\documentclass[11pt,a4paper]{article}

\usepackage[utf8]{inputenc}
\usepackage[T1]{fontenc}
\usepackage{lmodern}
\usepackage{geometry}
\usepackage{hyperref}
\usepackage{enumitem}

\geometry{margin=2.5cm}

\title{\textbf{GC-60 Model Overview}\\
\large Structural framework for prime number exploration}

\author{}
\date{}

\begin{document}

\maketitle

\section*{Introduction}

GC-60 is an arithmetic model designed to explore prime numbers by reducing the natural numbers to a structured and admissible subset.
The model is based on residue classes modulo 60 and is intended to support efficient exploration of large numeric ranges.

GC-60 is not an algorithm by itself, but a \textbf{conceptual framework} that can be implemented in different computational forms.

\section*{Motivation}

Classical approaches to prime number exploration often rely on exhaustive sieving or brute-force testing over contiguous numeric intervals.
Such methods become inefficient or impractical when the goal is to explore sparse or distant regions of the number line.

GC-60 addresses this limitation by:
\begin{itemize}[leftmargin=1.5cm]
\item reducing the search space using arithmetic constraints;
\item separating data construction from data interrogation;
\item enabling localized exploration within a globally defined structure.
\end{itemize}

\section*{Residue classes modulo 60}

Every integer can be expressed as:
\[
n = 60k + r, \quad r \in \{0,1,\dots,59\}.
\]

Among these residues, only a subset can contain prime numbers greater than 5.
GC-60 restricts attention to the following admissible residue classes:
\[
\{1,7,11,13,17,19,23,29,31,37,41,43,47,49,53,59\}.
\]

This reduction eliminates all integers divisible by 2, 3, or 5 at the structural level.

\section*{Archive concept}

GC-60 introduces the concept of an \textbf{arithmetic archive}.

The archive stores structured information derived from the admissible residue classes and supports:
\begin{itemize}[leftmargin=1.5cm]
\item incremental construction;
\item persistence across executions;
\item extension without invalidating previous data.
\end{itemize}

Each archive state defines a maximum exploration capacity, typically related to the square of the largest covered divisor.

\section*{Windows of exploration}

Instead of operating on a fixed numeric interval, GC-60 allows the definition of arbitrary \textbf{windows of exploration}.

A window is defined by:
\[
[\,A, B\,] \subset \mathbb{N}
\]
and can be explored provided that the archive contains sufficient information to cover the required divisors up to \(\sqrt{B}\).

This design enables localized exploration of very large numbers without global recomputation.

\section*{Separation of concerns}

A fundamental principle of GC-60 is the separation between:
\begin{itemize}[leftmargin=1.5cm]
\item \textbf{archive construction}, which is computationally intensive but performed incrementally;
\item \textbf{prime extraction}, which is fast and window-specific.
\end{itemize}

This separation allows different implementations and analysis tools to coexist over the same archive.

\section*{Scope and limitations}

GC-60 does not aim to replace classical sieve algorithms nor to provide a general factorization method.
Its purpose is exploratory and structural:
to offer a different way of navigating the space of prime numbers.

Formal mathematical analysis and performance guarantees are left to future work.

\section*{Conclusion}

GC-60 defines a structured arithmetic space in which prime numbers can be explored incrementally and selectively.
Its value lies in the combination of modular reduction, persistence, and separation between data and interpretation.

\end{document}
